%; whizzy paragraph -pdf xpdf -latex ./whizzypdfptex.sh
%; whizzy-paragraph "^\\\\begin{frame}\\|\\\\emtext"
% latex beamer presentation.
% platex, latex-beamer でコンパイルすることを想定。 

%     Tokyo Debian Meeting resources
%     Copyright (C) 2012 Junichi Uekawa

%     This program is free software; you can redistribute it and/or modify
%     it under the terms of the GNU General Public License as published by
%     the Free Software Foundation; either version 2 of the License, or
%     (at your option) any later version.

%     This program is distributed in the hope that it will be useful,
%     but WITHOUT ANY WARRANTY; without even the implied warreanty of
%     MERCHANTABILITY or FITNESS FOR A PARTICULAR PURPOSE.  See the
%     GNU General Public License for more details.

%     You should have received a copy of the GNU General Public License
%     along with this program; if not, write to the Free Software
%     Foundation, Inc., 51 Franklin St, Fifth Floor, Boston, MA  02110-1301 USA

\documentclass[cjk,dvipdfmx,12pt]{beamer}
\usetheme{Tokyo}
\usepackage{monthlypresentation}

%  preview (shell-command (concat "evince " (replace-regexp-in-string "tex$" "pdf"(buffer-file-name)) "&")) 
%  presentation (shell-command (concat "xpdf -fullscreen " (replace-regexp-in-string "tex$" "pdf"(buffer-file-name)) "&"))
%  presentation (shell-command (concat "evince " (replace-regexp-in-string "tex$" "pdf"(buffer-file-name)) "&"))

%http://www.naney.org/diki/dk/hyperref.html
%日本語EUC系環境の時
\AtBeginDvi{\special{pdf:tounicode EUC-UCS2}}
%シフトJIS系環境の時
%\AtBeginDvi{\special{pdf:tounicode 90ms-RKSJ-UCS2}}

\newenvironment{commandlinesmall}%
{\VerbatimEnvironment
  \begin{Sbox}\begin{minipage}{1.0\hsize}\begin{fontsize}{8}{8} \begin{BVerbatim}}%
{\end{BVerbatim}\end{fontsize}\end{minipage}\end{Sbox}
  \setlength{\fboxsep}{8pt}
% start on a new paragraph

\vspace{6pt}% skip before
\fcolorbox{dancerdarkblue}{dancerlightblue}{\TheSbox}

\vspace{6pt}% skip after
}
%end of commandlinesmall

\title{東京エリアDebian勉強会}
\subtitle{第116回 2014年8月度}
\author{野島貴英}
\date{2014年8月23日}
\logo{\includegraphics[width=8cm]{image200607/openlogo-light.eps}}

\begin{document}

\begin{frame}
\titlepage{}
\end{frame}

\begin{frame}{設営準備にご協力ください。}
会場設営よろしくおねがいします。
\end{frame}

\begin{frame}{Agenda}
 \begin{minipage}[t]{0.45\hsize}
  \begin{itemize}
   \item 注意事項
	 \begin{itemize}
	  \item 写真はセミナールーム内のみ可です。
          \item 出入りは自由でないので、もし外出したい方は、野島まで一声くださいませ。
	 \end{itemize}
   \item 事前課題発表
  \end{itemize}
 \end{minipage} 
 \begin{minipage}[t]{0.45\hsize}
  \begin{itemize}
   \item 最近あったDebian関連のイベント報告
	 \begin{itemize}
	  \item 第115回 東京エリアDebian勉強会 with 第2回 Debianパッケージング道場
	 \end{itemize}
   \item Debian Trivia Quiz
   \item Debianでタイルマップサービス作ってみた
   \item 今後のイベント
   \item 今日の宴会場所
  \end{itemize}
 \end{minipage}
\end{frame}

\section{事前課題}
\emtext{事前課題}
{\footnotesize
\begin{prework}{ dictoss(杉本　典充) }
\begin{itemize}
\item 前回のパッケージング道場の続きをする
\item VAIO Pでlinux-image-3.14の場合にMP3再生が倍速再生されるバグの調査
\end{itemize}
\end{prework}

\begin{prework}{ 吉野(yy\_{}y\_{}ja\_{}jp) }
DDTSS
\end{prework}

\begin{prework}{ zinrai }
iodine,nikolaを使ってみる。\\
\url{https://packages.debian.org/sid/iodine}\\
\url{https://packages.debian.org/sid/nikola}
\end{prework}

\begin{prework}{ 田上　健太 }
新しく個人用ノートPCを買ったので、開発環境を構築していきます。
\end{prework}

\begin{prework}{ 野島　貴英 }
xmrisパッケージ化やる。
\end{prework}

\begin{prework}{ wbcchsyn }
6月に発表したツールについて、頂いた意見の取り込みとフィードバックをおこないます。
\end{prework}

\begin{prework}{ まえだこうへい }
\url{http://d.palmtb.net/pages/about.html#todo-from-maintainer-dashboard}
のバグ潰し
\end{prework}

\begin{prework}{ 野首(@knok) }
\begin{itemize}
\item KAKASIのローマ字テーブル見直し
\item Leap Motion, HID送信機のハック
\item ECS Livaの問題調査
\end{itemize}
\end{prework}

}

\section{イベント報告}
\emtext{イベント報告}

\begin{frame}{第115回東京エリアDebian勉強会with第2回Debianパッケージング道場報告}
 
\begin{itemize}
\item  (株)サイバーエージェントさんをお借りしての開催でした。途中参加の方も含め参加者10名。
\item 10:00-16:00までは第2回Debianパッケージング道場を岩松さんにより、
\item 16:00以降はいつもの東京エリアDebian勉強会の形式で、前田さんにより「Jenkinsを
用いたパッケージ作成の自動化」の発表
\end{itemize} 

 両者とも最近のトレンドを踏まえた内容でした。
 
\end{frame}

\section{Debian Trivia Quiz}
\emtext{Debian Trivia Quiz}
\begin{frame}{Debian Trivia Quiz}

  Debian の常識、もちろん知ってますよね?
知らないなんて恥ずかしくて、知らないとは言えないあんなことやこんなこと、
みんなで確認してみましょう。

今回の出題範囲は\url{debian-devel-announce@lists.debian.org},
\url{debian-devel@lists.debian.org} に投稿された
内容などからです。

\end{frame}

\subsection{問題}

%; whizzy-master ../debianmeetingresume201311.tex
% 以上の設定をしているため、このファイルで M-x whizzytex すると、whizzytexが利用できます。
%

\santaku
{2014/8/16にDebianは誕生日を迎えました。さて今年で何歳?}
{25}
{21}
{20}
{B}
{1993/8/16にIan Ashley Murdockさんにより、Debianは創設されました。というわけで、21歳です。Debian Dayはもう過ぎちゃいましたが、本日(2014/8/23)の勉強会後の宴会で祝いたいと思います！}

\santaku
{2014/7時点の各国のアクティブなDebian Developerの数とそれぞれの国の人口との比率が最も多いのはどこの国?}
{Finland}
{Ireland}
{New Zealand}
{A}
{Finlandは本割合は2009年からずっと毎年1位(今年3.61ppm。)2位 Ireland 3位 New Zealand。日本の本割合は0.28ppmで30位でした。単にアクティブDeveloperの絶対数が多い国は1位USA、2位Germany、3位France。}

\santaku
{OpenAmbitが2014/7にパッケージ化されsidへリリースされました。ところで何するパッケージ?}
{AmazonTV用アプリ開発環境}
{ChoromeCast用アプリ開発環境}
{SUNNTO AMBIT用アプリ開発環境}
{C}
{SUNNTO AMBITという腕時計型のGPS搭載の端末のアプリ開発環境となります。楽天とか、Amazonで買えます。スポーツ好きな方はDebianでアプリ開発やってみてください。}

\santaku
{2014/7/31にtechnical committeeによりDebianのデフォルトのJPEG圧縮伸長ライブラリとして採択されたのは以下のうちのどれ?}
{libjpeg6b}
{libjpeg8/9}
{libjpeg-turbo}
{C}
{libjpeg-turboはlibjpeg6 ABI互換のJPEG圧縮伸長ライブラリで、libjpeg6を遥かに凌ぐ速さで動作します。ここで、JPEG圧縮伸長ライブラリのデフォルトとして、libjpeg8/9か、libjpeg-turboかで議論が分かれていましたが、technical comitteeはlibjpeg-turboをデフォルトのJPEG圧縮伸長ライブラリとしました。}

\santaku
{2014/7/31にtechnical committeeから、「Debianパッケージは複数のinitシステムに対応する事」ということが再アナウンスされました。何のパッケージのドタバタがきっかけでしょう?}
{ftp}
{tftp-hpa}
{ncftp}
{B}
{experimentalにtftp-hpa 5.2-17が入ったが、そのchangelogとして''Removing upstart hacks, they are ugly and upstart is dead now.''という内容が入っていた件が物議を醸しました。bug746715参照。}

\santaku
{2014/7/20にsqueeze（注：squeeze-ltsではない）の最後のアップデートが行われました。何回目のアプデートでしょうか?}
{10}
{9}
{8}
{A}
{最後のsqueezeのアップデートです。実はこのアップデートを行っても、Debianですでに報告されているsqueezeのセキュリティホールの一部は残ったままです。しかしながら、これ以上のアップデートは無いので、wheezyにアップデートするか、squeeze-ltsの採用を検討しましょう。}

\santaku
{2014/7/31にJessieに搭載の可能性のあるlinux kernelのバージョンのアナウンスが行われました。バージョンはいくつでしょう?}
{3.14}
{3.12}
{3.16}
{C}
{2014/8/14にkernel.orgでlinux kernel3.16のリリースがアナウンスされました。Debian Kernelチームのアナウンスでは、「3.16がJessieに採用される可能性があるので、現在のパッケージが3.16と互換のあるKernel APIの元で動作するか確認し、問題があれば然るべき対応してほしい」とのこと。}

\santaku
{VCS-*フィールドのVCS情報を参照し、パッケージのchangelogが合致しているかのチェックが行われるようになりました。以下のどのサイトで確認できる?}
{http://qa.debian.org/cgi-bin/vcswatch}
{http://bugs.debian.org/}
{http://codesearch.debian.net/}
{A}
{vcswatchのページにアクセスし、見たいパッケージ名入れると、状況が得られます。詳しくはvcswatchのサイトの説明をご覧ください。}

\santaku
{パッケージに含まれているドキュメントに関して、lintianが新たな警告をするようになりました。以下のどれ?}
{HTMLファイル中の画像/CSS/JS/videoリンクがローカルファイルを指していない場合、警告}
{HTMLファイル中のDebianの綴りが間違っている場合、警告}
{HTMLファイルが入っていない場合、警告}
{A}
{ブラウザがHTMLファイルを表示する場合、ローカルのファイルであっても、表示時にブラウザが自動でリンク先をアクセスしてしまう作用を持つリンクがあります。ここに外部サイトのURLが入っていた場合、その外部サイトは、HTMLファイルが閲覧された時の情報を収集することが出来てしまいます。こういったことはセキュリティ上良くないので、該当するようなHTMLファイルが含まれている場合、lintianが警告するようになりました。}

\santaku
{2014/7にBTSのWEBサイトにreplyのリンクがつくようになりました。このリンクをブラウザからクリックすると何が起きるようになった?}
{BTS上でreplyしている次のメッセージが表示されるようになった。}
{BTS上のreply用のフォームがブラウザに表示されるようになった。}
{メーラが起動し、正しいSubject/宛先/引用が入るようになった。}
{C}
{やってみれば判りますが、Bugレポートに対する返答をする場合に、非常に便利です。BTSから見て、正しい内容と宛先で、返答出来るように下書きメールがメーラに準備されるようになりました。}

\santaku
{2014/8/14にDebian Installer Jessie Beta 1がアナウンスされました。このバージョンのインストーラで導入されるinitシステムは結局どうなった?}
{systemdになった}
{sysvinitになった}
{upstartになった}
{A}
{このバージョンのDebian Installerから、initシステムとしてsystemdが入るようになりました。ちなみにこの件、Installerプログラム本体の話であって、Jessie本体のBeta版リリースという意味では無い事に注意。}


\section{Debianでタイルマップサービス作ってみた}
\emtext{Debianでタイルマップサービス作ってみた}

\section{今後のイベント}
\emtext{今後のイベント}
\begin{frame}{今後のイベント}
\begin{itemize}
 \item 関西エリアDebian勉強会
 \item 2014年9月 東京エリアDebian勉強会
\end{itemize}
\end{frame}

\section{今日の宴会場所}
\emtext{今日の宴会場所}
\begin{frame}{今日の宴会場所}
未定
\end{frame}

\end{document}

;;; Local Variables: ***
;;; outline-regexp: "\\([ 	]*\\\\\\(documentstyle\\|documentclass\\|emtext\\|section\\|begin{frame}\\)\\*?[ 	]*[[{]\\|[]+\\)" ***
;;; End: ***
